\documentclass[a4paper, 11pt]{article}
% ── Packages ──────────────────────────────────────────────────────────────────
\usepackage[utf8]{inputenc}
\usepackage[T1]{fontenc}
\usepackage[french]{babel}
\usepackage{amsmath, amssymb, amsfonts}
\usepackage{graphicx}
\usepackage{booktabs}
\usepackage{geometry}
\usepackage{hyperref}
\usepackage{enumitem}
\usepackage{caption}
\usepackage{subcaption}
\usepackage{float}
\geometry{margin=2.5cm}
\hypersetup{colorlinks=true, linkcolor=blue!60!black, citecolor=blue!60!black, urlcolor=blue!60!black}
\title{\textbf{Note d'étape 2 -- StatApp}\\
\large Incertitude des Large Language Models}
\author{ENSAE Paris -- Société Générale}
\date{16 mars 2026}
\begin{document}
\maketitle
% =============================================================================
\section{Rappel du contexte et objectifs}
% =============================================================================
Les Large Language Models (LLMs) génèrent des réponses textuelles fluides mais sans indication de leur degré de confiance. Notre projet, encadré par la Société Générale, vise à implémenter et évaluer différentes méthodes de quantification de l'incertitude des LLMs, en distinguant les approches \emph{white-box} (accès aux probabilités internes) des approches \emph{black-box} (sans accès aux logits). Depuis la première note d'étape (décembre 2025), nous avons considérablement élargi le périmètre du projet : nouveaux datasets, nouvelles méthodes, étude du lien entre confiance et proximité à la vérité, et premiers résultats comparatifs.
% =============================================================================
\section{Travail réalisé depuis la note d'étape 1}
% =============================================================================
\subsection{Extension des datasets}
Nous avons enrichi notre base d'évaluation au-delà de MMLU et GSM8K avec trois nouveaux datasets :
\begin{itemize}[nosep]
    \item \textbf{ELI5} (\emph{Explain Like I'm 5}) : questions ouvertes nécessitant des réponses longues explicatives. Ce dataset permet d'évaluer l'incertitude sur des générations en texte libre, contrairement aux QCM.
    \item \textbf{TriviaQA} : questions factuelles avec réponses courtes. Utilisé d'une part pour mesurer l'importance des tokens, et d'autre part, en filtrant les questions à réponse numérique, pour l'étude du lien entre confiance et proximité à la vérité.

\end{itemize}
\subsection{Nouvelles méthodes implémentées}
Au-delà des méthodes white-box (log-probabilités) et SelfCheckGPT implémentées lors de la première phase, nous avons développé quatre nouvelles approches :
\paragraph{Entropie sémantique (Kuhn et al., 2023).}
Pour chaque question, nous générons $K=5$ réponses avec température élevée ($T=1{,}0$) et log-probabilités. Un modèle NLI (Gemini~2.5 Flash) regroupe les réponses sémantiquement équivalentes via un algorithme Union-Find. L'entropie sémantique est ensuite calculée comme l'entropie de Shannon sur les probabilités des clusters :
\begin{equation}
    \text{SE} = -\sum_{i} P(C_i) \log P(C_i), \quad P(C_i) = \frac{\sum_{r \in C_i} \exp(\bar{\ell}_r)}{\sum_{j} \sum_{r \in C_{j}} \exp(\bar{\ell}_r)}
\end{equation}
où $\bar{\ell}_r$ désigne la log-probabilité moyenne (normalisée par la longueur) de la réponse $r$. Une entropie sémantique nulle indique que toutes les réponses sont équivalentes (haute confiance), tandis qu'une valeur élevée signale des réponses sémantiquement divergentes.
\paragraph{Prédiction conforme (LAC et APS).}
Nous avons implémenté deux méthodes de prédiction conforme sur MMLU : \emph{Least Ambiguous set-valued Classifier} (LAC) et \emph{Adaptive Prediction Sets} (APS). Pour un taux d'erreur nominal $\alpha$, ces méthodes construisent des ensembles de prédiction $C(X_t)$ garantissant $P(Y_t \in C(X_t)) \geq 1-\alpha$. Le protocole divise les données en deux (50\% calibration, 50\% test), calcule un quantile $\hat{q}$ sur les scores de non-conformité, puis construit les ensembles sur le jeu de test.
\paragraph{Importance des tokens par perplexité et similarité cosinus.}
Sur TriviaQA, nous avons développé deux méthodes d'attribution d'importance aux tokens :
\begin{itemize}[nosep]
    \item \emph{Perplexity importance} : $\text{imp}(t) = -\log p(t)$, les tokens de faible probabilité sont les plus informatifs.
    \item \emph{Cosine importance} : $\text{imp}(t) = 1 - \cos(\vec{v}_{\text{full}}, \vec{v}_{\setminus t})$, où $\vec{v}_{\text{full}}$ est le vecteur moyen des distributions top-$k$ sur toutes les positions et $\vec{v}_{\setminus t}$ le même vecteur calculé en excluant la position du token $t$.
\end{itemize}
Ces méthodes sont évaluées aux niveaux token et mot, avec comparaison aux vérités du dataset.
\paragraph{Lien confiance--proximité à la vérité.}
Sur un sous-ensemble de TriviaQA filtré aux questions à réponse numérique ($n=3\,395$), nous avons étudié la relation entre différents scores de confiance (verbalisée, jugement externe, log-probabilités) et la distance à la vérité terrain, mesurée par erreur absolue, relative et logarithmique (cf.\ section~\ref{sec:confiance_verite}).
% =============================================================================
\section{Méthodologie}
% =============================================================================
\subsection{Cadre expérimental commun}
Toutes les expériences textuelles utilisent Gemini~2.0 Flash via l'API Vertex AI avec \texttt{response\_logprobs=True}. Pour les évaluations nécessitant un juge LLM (SelfCheckGPT, évaluation de justesse), nous utilisons Gemini~2.5 Flash. Le code est structuré en modules Python avec résultats sauvegardés en JSONL pour permettre la reprise après interruption.
\subsection{Méthodes white-box : signaux de confiance}
Nous comparons plusieurs signaux extraits des log-probabilités :
\begin{itemize}[nosep]
    \item \textbf{Entropie normalisée} ($c_\text{ent}$) : $H_\text{norm} = -\frac{1}{\log K}\sum_{k} p_k \log p_k$ sur les probabilités des options A/B/C/D. Confiance $= 1 - H_\text{norm}$.
    \item \textbf{Probabilité maximale} ($c_\text{pmax}$) : $\max_k p_k$, la probabilité de l'option la plus probable.
    \item \textbf{Confiance déclarée} ($c_\text{decl}$) : le modèle verbalise lui-même un score de confiance entre 0 et 1.
    \item \textbf{Perplexité de séquence} : $\text{PPL} = \exp\left(-\frac{1}{N}\sum_i \log p(t_i | t_{<i})\right)$ sur la réponse complète.
    \item \textbf{Entropie sémantique} : mesure la diversité sémantique des réponses échantillonnées (cf. section précédente).
\end{itemize}
\subsection{Méthodes black-box}
\begin{itemize}[nosep]
    \item \textbf{SelfCheckGPT} : pour chaque phrase de la réponse, un juge LLM vérifie sa cohérence avec $K$ passages alternatifs générés pour la même question. Score $\in [0,1]$ : 0 = fiable, 1 = contredit.
\end{itemize}
\subsection{Évaluation des méthodes}
Pour comparer les signaux de confiance, nous utilisons des courbes \emph{accuracy vs coverage} : en ne retenant que les réponses les plus confiantes, l'accuracy devrait augmenter si le signal est discriminant. Nous mesurons également : le score de Brier, l'ECE (\emph{Expected Calibration Error}), l'AUROC et l'AUPRC. La prédiction conforme est évaluée par sa couverture empirique et la taille moyenne des ensembles de prédiction.
% =============================================================================
\section{Statistiques descriptives}
% =============================================================================
\subsection{MMLU}
Les statistiques descriptives de MMLU (15\,858 questions, 57 sujets) ont été présentées dans la note d'étape 1. Pour l'évaluation comparative des méthodes d'incertitude, nous travaillons sur un sous-ensemble de 107 questions disposant de logits complets A/B/C/D et de confiance déclarée, permettant une comparaison appariée stricte.
\subsection{ELI5}
100 questions ELI5 traitées par Gemini~2.0 Flash produisent des réponses de 3 à 5~phrases. La perplexité médiane par phrase est de 1,31 (distribution asymétrique à droite), tandis que l'entropie top-$k$ médiane est de 0,52. L'entropie sémantique, calculée à partir de $K=5$ réponses par question, présente une distribution fortement bimodale (figure~\ref{fig:se_distrib}) : 51 questions ont une SE nulle (toutes les réponses sémantiquement équivalentes), tandis que les autres se répartissent entre 0,4 et 1,6. Un juge LLM attribue un score de qualité $\geq 4$ à 82\% des réponses.

\begin{figure}[H]
    \centering
    \includegraphics[width=0.55\textwidth]{figures/semantic_entropy_distribution.png}
    \caption{Distribution de l'entropie sémantique sur 100 questions ELI5.}
    \label{fig:se_distrib}
\end{figure}

\subsection{TriviaQA}
100 questions TriviaQA avec réponses courtes. Le modèle atteint un taux de justesse de 87\% selon le juge LLM. L'analyse porte d'une part sur la comparaison entre les tokens identifiés comme importants par nos méthodes et les tokens de la réponse correcte, et d'autre part, pour l'étude confiance--proximité (section~\ref{sec:confiance_verite}), sur un sous-ensemble de 3\,395 questions à réponse numérique, permettant de définir des distances interprétables entre la réponse du modèle et la vérité terrain.
% =============================================================================
\section{Premiers résultats}
% =============================================================================
\subsection{Comparaison des signaux de confiance sur MMLU}
Sur les 107 questions appariées, nous comparons les courbes accuracy-coverage des différents signaux (figure~\ref{fig:acc_cov}). Les signaux basés sur les log-probabilités surpassent nettement la confiance déclarée : en sélectionnant les 10\% de questions les plus confiantes, l'accuracy atteint 81,8\% avec les logprobs contre 72,7\% avec la confiance déclarée. Pour les 50\% les plus confiantes, l'écart se creuse davantage (75,9\% vs 51,9\%). La figure~\ref{fig:distributions} confirme ce résultat : les distributions de confiance basées sur les logprobs séparent bien les réponses correctes des incorrectes, tandis que les distributions déclarées se chevauchent fortement.

\begin{figure}[H]
    \centering
    \includegraphics[width=0.70\textwidth]{figures/fig2_accuracy_coverage_paired.png}
    \caption{Courbes accuracy vs couverture sur MMLU (107 questions appariées).}
    \label{fig:acc_cov}
\end{figure}

\begin{figure}[H]
    \centering
    \includegraphics[width=0.85\textwidth]{figures/fig5_confidence_distributions.png}
    \caption{Distributions des scores de confiance pour réponses correctes vs incorrectes.}
    \label{fig:distributions}
\end{figure}

\paragraph{Prédiction conforme.}
Les garanties de couverture sont respectées empiriquement (figure~\ref{fig:conformal}). Pour $\alpha = 0{,}05$, LAC atteint une couverture de 96,2\% (cible : 95\%) avec une taille moyenne d'ensemble de 3,2, tandis qu'APS atteint 100\% avec des ensembles de taille 4,0. APS tend à être plus conservateur mais fournit des garanties plus stables.

\begin{figure}[H]
    \centering
    \includegraphics[width=\textwidth]{figures/fig3_conformal_prediction.png}
    \caption{Prédiction conforme sur MMLU : couverture et taille d'ensemble vs $\alpha$.}
    \label{fig:conformal}
\end{figure}

\subsection{Incertitude sur les réponses longues (ELI5)}
L'analyse sur ELI5 révèle que l'entropie sémantique est anti-corrélée avec la qualité des réponses (Spearman $r = -0{,}173$, $p = 0{,}08$) : les questions à SE nulle obtiennent les meilleurs scores d'accuracy. La figure~\ref{fig:eli5_coverage} montre les courbes accuracy-coverage : le signal combiné (moyenne des confiances normalisées) offre un bon compromis, restant au-dessus de la baseline pour toutes les couvertures.

\begin{figure}[H]
    \centering
    \includegraphics[width=0.85\textwidth]{figures/eli5_accuracy_coverage.png}
    \caption{Courbes accuracy vs couverture sur ELI5.}
    \label{fig:eli5_coverage}
\end{figure}

\subsection{Importance des tokens sur TriviaQA}
La figure~\ref{fig:hit_rate} compare les \emph{hit rates} des méthodes d'importance. L'analyse au niveau mot surpasse systématiquement celle au niveau token, et la similarité cosinus domine la perplexité. Un \emph{hit rate} autour de 85\% pour la similarité cosinus montre que cette méthode est efficace pour localiser l'information clé au sein d'une phrase.

\begin{figure}[H]
    \centering
    \includegraphics[width=0.55\textwidth]{figures/judge_hit_rate.png}
    \caption{Hit rates des méthodes d'importance sur TriviaQA.}
    \label{fig:hit_rate}
\end{figure}

\subsection{Lien entre confiance et proximité à la vérité}
\label{sec:confiance_verite}
Cette analyse, menée sur le sous-ensemble de TriviaQA à réponses numériques ($n=3\,395$, accuracy $\approx 86\%$), étudie si un score de confiance élevé implique une réponse plus proche de la vérité terrain. La distance est mesurée par l'erreur absolue $|\hat{y}_i - y_i|$, l'erreur relative $|\hat{y}_i - y_i|/(|y_i|+\varepsilon)$ et l'erreur logarithmique $|\log \hat{y}_i - \log y_i|$.

Plusieurs scores de confiance sont comparés : la \emph{confiance verbalisée} (le modèle annonce lui-même un score $\in [0,1]$), la \emph{confiance par jugement externe} (un second LLM évalue la plausibilité), et des scores fondés sur les log-probabilités internes (probabilité jointe, moyenne géométrique, perplexité, probabilité minimale), calculés soit sur l'ensemble de la génération, soit restreints aux tokens de la réponse numérique.

\paragraph{Discrimination et calibration.}
Les scores fondés sur les log-probabilités des tokens de la réponse obtiennent les meilleures AUC (0,84--0,86), devant la confiance verbalisée et le jugement externe (AUC $\approx 0{,}68$--$0{,}72$), et loin devant les métriques calculées sur l'ensemble de la génération (AUC $\approx 0{,}60$--$0{,}63$). Tous les scores présentent une surconfiance systématique : les moyennes de confiance (0,91--0,99) dépassent l'accuracy empirique (0,86), avec des scores de Brier entre 0,11 et 0,21.

\paragraph{Corrélation confiance--distance.}
Sur l'ensemble des réponses, la corrélation de Spearman entre confiance et distance à la vérité est négative (les réponses plus confiantes sont plus proches de la vérité). Les corrélations les plus fortes proviennent des log-probabilités restreintes aux tokens de la réponse (Spearman $\approx -0{,}49$), suivies de la confiance verbalisée et du jugement externe ($\approx -0{,}27$). Cependant, lorsqu'on restreint l'analyse aux seules réponses incorrectes ($n=463$), les corrélations chutent fortement (Spearman $\approx -0{,}13$ à $-0{,}22$ pour les meilleurs signaux), indiquant que la confiance permet surtout de séparer réponses correctes et incorrectes, mais quantifie mal l'amplitude de l'erreur parmi les réponses fausses. Les résultats sont robustes au choix de la mesure de distance.

% =============================================================================
\section{Perspectives pour le mémoire final}
% =============================================================================
La suite du projet s'articulera autour de trois axes :
\begin{enumerate}[nosep]
    \item \textbf{Consolidation des résultats quantitatifs.} Augmenter le nombre de questions traitées sur ELI5 et TriviaQA, produire des tables de métriques complètes pour chaque combinaison méthode$\times$dataset, et effectuer des tests de significativité statistique.
    \item \textbf{Construction d'un score d'incertitude composite.} Proposer une méthode de fusion des différents signaux (white-box et black-box) pour construire un score global, et évaluer sa performance sur les différents datasets.
    \item \textbf{Synthèse comparative et recommandations.} Formaliser une comparaison systématique white-box vs black-box selon les axes coût computationnel, qualité de calibration et applicabilité pratique.
\end{enumerate}
\vfill
\section*{Annexes}
Les résumés des articles, fiches méthodologiques, code source, figures des statistiques descriptives et résultats intermédiaires sont disponibles sur le Github du projet.
\url{https://github.com/gregouzeee/statapp}
\end{document}
